\rok{Септембар 2025.}{K}

Први практични тест из Алгоритама и структура података

\zadatak{1}{Quick Sort}
Написати \textit{quick} методу \textit{Quick Sort} алгоритма за сортирање целобројног низа дужине \textit{n}.

\textbf{Додатне напомене за решавање задатка:}
\begin{itemize}
    \item Улаз је несортиран низ целих бројева.
    \item Излаз је сортирани низ целих бројева.
    \item У шаблону се налазе класе \texttt{RandomArrayGenerator} и \texttt{RandomNumberGen} које треба искористити за генерисање низа случајних бројева.
    \item У оквиру класе \texttt{QuickSort} написати имплементацију за методу:
    \begin{Verbatim}[fontsize=\footnotesize, numbers=left, xleftmargin=10mm]
private static void quickSort(int[] arr, int lower, int upper)
    \end{Verbatim}
    \item Обезбедити код у \texttt{Main} методи за тестирање алгоритма.
\end{itemize}



\zadatak{2}{Промена редоследа последњих n елемената стека користећи ред}
Имплементирати функцију која за задати стек целих бројева обрће редослед само последњих n елемената стека, користећи помоћни ред. У оквиру тих последњих n елемената, за сваки од њих који је непаран важиће да се његова вредност умањи за 2, а за сваки од тих n елемената који је паран његова вредност ће се удвостручити. Уколико након умањивања непарног броја за 2 добијемо негативан број, вредност тог броја се поставља на 1. Осим тога, остали елементи стека остају у истом редоследу.

\textbf{Спецификација функције:}
\begin{itemize}
    \item Функција прима број n, где је n број елемената које треба обрнути.
    \item Функција треба да користи ред (\textit{queue}) као помоћну структуру.
    \item Елементи пре прослеђених n треба да остану у оригиналном редоследу.
\end{itemize}

\textbf{Додатни захтеви за решавање задатка:}
\begin{itemize}
    \item Алгоритам писати у оквиру класе \texttt{Stack}.
    \item Захтеван потпис методе:
    \begin{Verbatim}[fontsize=\footnotesize, numbers=left, xleftmargin=10mm]
public void ReverseAndModifyLastNElements (int n)
    \end{Verbatim}
\end{itemize}

\textbf{Пример 1:}
\begin{itemize}
    \item \textbf{Оригинални стек:} 4 3 1 10 9 4, n = 3 \\
    Међу последња 3 елемента, 10 и 4 су парани бројеви чија вредност ће се удвостручити, а број 9 је непаран и његова вредност се смањује за 2.
    \item \textbf{Стек након обртања 3 елемента:} 4 3 1 8 7 20
\end{itemize}

\textbf{Пример 2:}
\begin{itemize}
    \item \textbf{Оригинални стек:} 4 3 1 1 2 3, n = 3
    \item \textbf{Стек након обртања последњих 3 елемената:} 4 3 1 1 4 1
\end{itemize}

\textbf{Пример 3:}
\begin{itemize}
    \item \textbf{Оригинални стек:} 4 3 1 4 3 2, n = 8
    \item Неважећа вредност за n.
\end{itemize}



\zadatak{3}{Број и производ „увала" и „брегова" у једноструко повезаној листи}
Уједноструко спрегнутој листи сваки од чворова може да представља увалу или брег. Чвор ће представљати \textbf{увалу} уколико је његова вредност мања од оба суседна елемента (елемент пре и после посматраног чвора у оквиру листе) – додатно, разлика у односу на оба суседа мора бити најмање 2. Чвор ће представљати \textbf{брег} уколико је његова вредност већа од оба суседна елемента (елемент пре и после посматраног чвора у оквиру листе) – додатно, разлика у односу на оба суседа мора бити најмање 3. Метода враћа низ у коме прва вредност представља број увала и брегова, а друга производ њихових вредности. Напомена: чвор може бити увала или брег само ако има претходни и следећи чвор.

\textbf{Додатни захтеви за решавање задатка:}
\begin{itemize}
    \item Решавај овај задатак помоћу структуре листе (\texttt{LinkedList}) из шаблона задатка.
    \item Имплементирај га у оквиру методе:
    \begin{Verbatim}[fontsize=\footnotesize, numbers=left, xleftmargin=10mm]
public int[] CountAndMultiplyValleysAndPeaks()
    \end{Verbatim}
\end{itemize}

\textbf{Пример 1:}
\begin{itemize}
    \item \textbf{Улаз:} листа: $1 \to 3 \to 1 \to 4 \to 2 \to 7 \to 3 \to 6$
    \item \textbf{Излаз:} [4, 42]
    \item \textbf{Објашњење:} \\
    1 је мање за 2 или више од суседних 3 и 4. \\
    2 је мање за 2 или више од суседних 4 и 7. \\
    3 је мање за 2 или више од суседних 7 и 6. \\
    7 је веће за 3 или више од суседних 2 и 3. \\
    Увале: 1, 2 и 3. Брегови: 7. Укупан број је 4, а производ $1 \cdot 2 \cdot 3 \cdot 7$ је 42.
\end{itemize}

\textbf{Пример 2:}
\begin{itemize}
    \item \textbf{Улаз:} листа: $1 \to 2 \to 2 \to 5 \to 5$
    \item \textbf{Излаз:} [-1, -1]
    \item \textbf{Објашњење:} Немамо критичне тачке, односно увале и брегове јер ниједан од чворова не задовољава постављене услове.
\end{itemize}

\sablon{Шаблон првог теста новембар 2024. група K}{2024/Novembar/GrupaK/AISP_2025_Test1_GrupaK.linq}
